% ------------------------------- Chapter 3 -------------------------------
\chapter{Rotor Blade Design}

The rotor-blade design converts the mean-line flow angles into an actual blade
profile. A few geometric assumptions (aspect ratio, stagger) fix the chord and
the number of blades, after which the camber line and the thickness distribution
are built with ParaBlade from the inlet/exit metal angles.

% =====================================================================
\section{Input Data Assumptions}

\paragraph{Aspect ratio $AR = b/\bar c$.}
The aspect ratio relates the blade span $b$ to the chord $\bar c$. For
high-pressure rotors $AR$ lies between $1$ and $1.5$; a value $AR = 1$ is adopted,
which gives the best compromise for the present loading.

\paragraph{Stagger angle $\xi$.}
The stagger is the angle between the chord and the axial direction. It is read
from the Kacker \& Okapuu correlation~\cite{kacker} as a function of the inlet and
exit relative angles $\beta_2 = 42.64\degr$ and $\beta_3 = 69.53\degr$, giving
$\xi \approx 39\degr$.

\paragraph{Zweifel loading coefficient.}
The blade pitch (and hence the blade count) is set from the Zweifel
criterion~\cite{zweifel} with $Z = 1.0$, which fixes the optimum
pitch-to-axial-chord ratio for the prescribed turning.

% =====================================================================
\section{Design Procedure and Calculations}

From the aspect ratio and the span, the chord and axial chord are obtained; the
Zweifel criterion then fixes the pitch and, with the mean circumference, the
number of blades:
\begin{equation}
    N_b = \frac{2\pi R_m}{s},\qquad s = \text{pitch}.
\end{equation}
The camber line is generated between the metal angles $\beta_2$ (inlet) and
$\beta_3$ (exit) and combined with a prescribed thickness distribution to produce
the suction and pressure surfaces using ParaBlade~\cite{parablade}. The tool
reproduces the prescribed analytic curvature: a check of the curvature design
variable against the analytic surface curvature matches at both ends
($\kappa_{LE}=17.61$, $\kappa_{TE}=213.29$), confirming a consistent profile.

% =====================================================================
\section{Main Blade Parameters}

The resulting blade parameters are collected in Table~\ref{tab:blade}; the
constructed section is shown in Fig.~\ref{fig:section}.





\begin{minipage}{0.58\textwidth}
    \begin{table}[H]
\centering
\caption{Main rotor-blade parameters.}
\label{tab:blade}
\begin{tabular}{l r}
\toprule
\textbf{Parameter} & \textbf{Value}\\
\midrule
Inlet flow angle $\alpha_2$ [\degr]          & 71.10 \\
Inlet metal (rel.) angle $\beta_2$ [\degr]   & 42.64 \\
Exit flow angle $\alpha_3$ [\degr]           & 34.18 \\
Exit metal (rel.) angle $\beta_3$ [\degr]    & 69.53 \\
Aspect ratio $AR$                            & 1.00 \\
Stagger angle $\xi$ [\degr]                   & 39.0 \\
Zweifel coefficient $Z$                      & 1.00 \\
Chord $c$ [m]                                & 0.192 \\
Axial chord $c_x$ [m]                        & 0.149 \\
Pitch $s$ [m]                                & 0.170 \\
Pitch-to-axial-chord $s/c_x$                 & 1.14 \\
Number of blades $N_b$                       & 42 \\
Max. thickness-to-chord $t_{max}/c$          & 0.20 \\
Leading-edge radius $R_{LE}$ [m]             & 0.0085 \\
Trailing-edge radius $R_{TE}$ [m]            & 0.0007 \\
End-wall flare angle $\varepsilon$ [\degr]    & 15.90 \\
\bottomrule
\end{tabular}
\end{table}
\end{minipage}
\hfill
\begin{minipage}{0.38\textwidth}
    \begin{figure}[H]
\centering
\includegraphics[width=0.55\textwidth]{Blade Section.png}
\caption{Constructed rotor-blade section (camber line and thickness distribution).}
\label{fig:section}
\end{figure}
\end{minipage}